\section{Problem Statement}
\label{sec:problem-statement}

Given a pretrained, frozen visual encoder applied to the frames of a procedural video
clip, does the resulting latent trajectory $z_1, \dots, z_T$ already support two
things: (i) recovering \emph{how far the clip has progressed} through the task
(temporal ordinality), and (ii) recovering \emph{unseen intermediate states} given only
a start frame $z_1$ and a goal frame $z_T$ (planning by interpolation) --- without ever
training on ground-truth actions? And, separately, can a self-supervised training
objective induce this structure in an encoder that does not already have it, in a way
that transfers to held-out data rather than only fitting the training clips?


\subsection{Motivation}
Long-horizon procedural video understanding needs a notion of \emph{progress} --- where
along a task a clip currently is --- and the ability to reason about intermediate states
between a start and a goal, without access to ground-truth actions. Foundation visual
encoders (DINOv2, DINOv3, V-JEPA2, \ldots) are usually evaluated on semantic tasks
(classification, retrieval), not on whether their latent trajectories are temporally
ordered or ``straight'' enough to support interpolation-based planning. Meanwhile, many
self-supervised objectives claim to induce temporally-ordered, progress-aware
representations (progress regression, ordinal classification, pairwise ranking,
temporal contrastive losses, goal-conditioned value learning, \ldots), but these claims
are rarely tested against a common, encoder-agnostic evaluation protocol -- so it is
hard to tell whether a reported gain reflects the objective, the encoder, the metric, or
representation collapse.

\subsection{Goals}
\begin{itemize}
    \item Define scale-invariant, comparable ordinality and planning metrics that work
    across encoders with very different raw latent scales (progress-metrics:
    time-progress Spearman, straightness; interpolation planning: cosine similarity,
    relative $L_2$, $R^2$).
    \item Benchmark existing frozen encoders (DINOv2 variants, DINOv3, V-JEPA2 / 2.1) on
    these metrics with nothing trained, to see which latent spaces already support
    planning ``for free.''
    \item Train and score a wide set of candidate self-supervised objectives (absolute
    progress regression and its ordinal variants, pairwise/rank-based losses, temporal
    contrastive and alignment losses, goal-conditioned value learning) under one
    held-fixed encoder / optimizer / data recipe, so that any difference in the result is
    attributable to the objective alone.
    \item Build a predictive world model (encoder + predictor trained end-to-end on
    next-latent prediction, with optional curvature / VICReg regularization) as a
    complementary approach that does train the encoder itself on procedural video, rather
    than only probing a frozen one.
    \item Determine whether latent-space ``straightness'' is necessary and/or sufficient
    for planning, and whether apparent gains on ordinality metrics are confounded with
    representation collapse.
\end{itemize}
